\chapter{Biểu diễn số nhị phân}
\label{Chapter2}
\setcounter{secnumdepth}{5}

\section{}
\section{}
\section{}
\subsection{}
\subsection{}
\subsection{Chia dữ liệu train/test theo phân tầng}
Hàm \texttt{train\_test\_split\_stratified} được thiết kế để chia dữ liệu thành tập huấn luyện (train) và tập kiểm tra (test), đảm bảo rằng phân phối của giá trị mục tiêu (\texttt{Price}) trong hai tập này tương tự nhau. Điều này rất quan trọng vì nếu phân phối không đồng đều, mô hình có thể bị thiên lệch và không phản ánh đúng hiệu suất thực tế. Ví dụ, nếu tập train chứa chủ yếu các xe giá thấp và tập test chứa nhiều xe giá cao, mô hình sẽ không học được cách dự đoán giá xe cao cấp, dẫn đến sai số lớn.

Các bước thực hiện của hàm bao gồm:

\begin{itemize}
    \item \textbf{Bước 1: Chuẩn bị dữ liệu}: Giá trị mục tiêu \texttt{Price} đã được biến đổi logarit trước đó (sẽ được giải thích ở phần sau). Để chia dữ liệu theo phân tầng, chúng em chuyển giá trị này từ thang logarit về thang thực tế bằng hàm \texttt{np.expm1(y)}, trong đó \texttt{np.expm1(y)} tính toán $e^y - 1$, là phép nghịch đảo của \texttt{np.log1p}. Ví dụ, nếu giá trị logarit của \texttt{Price} là 13.8155, thì giá trị thực tế sẽ là $e^{13.8155} - 1 \approx 1,000,000$.
    \begin{figure}[h]
        \centering
        \includegraphics[scale=.5]{images/Price distribution.png}
        \caption{Phân bố của Price trước và sau khi sử dụng \texttt{np.log1p}}
    \end{figure}
\end{itemize}
\section{}
\section{}
\section{}
\subsection{Mô hình 2 - Chung Tín Đạt}
\subsubsection{Thêm đặc trưng bậc cao và tương tác}
Bổ sung thêm đặc trưng Area và $\text{Area}^2$ với công thức sau:
\begin{itemize}
    \item Area = Length $\times$ Width: Diện tích sàn xe.
    \item $\text{Area}^2$: Bình phương diện tích sàn xe.
\end{itemize}

\subsubsection{Các đặc trưng được sử dụng để huấn luyện}
\begin{lstlisting}
    features = ['Year', 'Max Power', 'Area^2', 'Make']
\end{lstlisting}

\subsubsection{Ý tưởng chọn đặc trưng}
\begin{enumerate}
    \item Về thực tế: Người dùng cơ bản thường có xu hướng mua những chiếc xe thuộc dòng mới nhất của 1 hãng nhất định. Do đó các hãng thường có chiến lược đẩy giá những xe mẫu mới so với những mẫu cũ cùng dòng, thể hiện ở đặc trưng Year. Bên cạnh thời điểm ra mắt, người dùng còn quan tâm đến việc xe chạy có ổn định, thoải mái hay không, điều đó được quyết định ở đặc trưng Max Power và Area. Ngoài ra, việc có nhiều hãng xe cạnh tranh với nhau trực tiếp ảnh hưởng đến thị trường. Các hãng thường tập trung vào 1 bộ phận khách hàng nhất định, tạo nên 1 phân khúc sản phẩm phù hợp với nhóm khách hàng đó, thể hiện ở đặc trưng Make. Tuy nhiên, vẫn có một số hãng có phân khúc rộng, còn chịu ảnh hưởng của từng dòng xe, hay Model. Nhưng để giảm sự phức tạp của mô hình, em đã không sử dụng đặc trưng Model.
    \item Về thống kê: Dựa vào bảng heatmap tương quan dưới đây (Bảng thể hiện mức độ thay đổi của Price khi một đặc trưng sơ cấp thay đổi):
    
    \includegraphics[scale=0.7]{images/Features heatmap.png}
    
    Có thể thấy Price có tương quan tương đối cao với Price, Engine, Max Power, Max Torque, Length, Width, Fuel Tank Capacity, Age, Area.
    Với nhóm động cơ, vì Engine, Max Power, Max Torque có tương quan rất cao (0.9) nên em chỉ chọn ra một đặc trưng là Max Power.
    Với nhóm kích thước xe, vì Length và Width có tương quan cao, có thể chọn một trong hai nhưng như đã giải thích ở trên, người dùng thường quan tâm đến diện tích hơn là kích thước của 1 chiều, do đó chọn đặc trưng Area.
    Với nhóm thời gian, Year và Age đều có thể được chọn, em quyết định chọn Year.
\end{enumerate}

\subsubsection{Một số thay đổi có thể xem xét}
\begin{enumerate}
    \item Một số đặc trưng phân loại như hệ thống truyền động Transmission, loại nhiên liệu Fuel Type có thể thay đổi hiệu suất của xe, trực tiếp ảnh hưởng Price.
    \item Địa điểm Location cùng màu xe Color cũng có thể tác động về mặt tâm lý, một số màu được ưa chuộng hơn ở các mẫu xe cao cấp nhưng tương đối kén với các mẫu thông dụng.
    \item Đặc trưng Area có thể được sử dụng nhưng $\text{Area}^2$ giúp mô hình hội tụ nhanh hơn.
\end{enumerate}

\subsubsection{Mô hình hồi quy tuyến tính}

Phương trình hồi quy tuyến tính:
\begin{center}
$\hat{y} = \omega_1 \times \text{Year} + \omega_2 \times \text{Max Power} + \omega_3 \times \text{Area}^2 + \omega \times \text{Make} + \beta$
\end{center}

Trong đó:
\begin{itemize}
    \item $\hat{y}$: Giá trị dự đoán ở thang chuẩn hóa (tương ứng $y_{\text{scaled}}$).
    \item $\omega_1, \omega_2, \dots, \omega_n$: Trọng số của từng đặc trưng.
    \item $\beta$: Hệ số chặn.
    \item Make: Tập hợp các cột nhận giá trị nhị phân từ one-hot encoding của biến phân loại.
\end{itemize}

\subsubsection{Giải thích cụ thể}
Các trọng số của từng đặc trưng sẽ được cập nhật qua mỗi lần chạy hết tập dữ liệu bằng thuật toán \textit{Gradient Descent}:
\begin{itemize}
    \item Hàm mất mát (Loss Function): Trung bình bình phương bé nhất (Mean Squared Error - MSE):
    \begin{center}
    $L=\frac{1}{n}*\sum^n_{i=1}(\hat{y}_i-y_i)^2$
    \end{center}
    
    Với $\hat{y}_i = \omega_1 \times \text{Year}_i + \omega_2 \times \text{Max Power}_i + \omega_3 \times \text{Area}^2_i + \omega_{ji} \times \text{Make}_{ji} + \beta$
    \item Mục tiêu: Tìm bộ trọng số $\omega, \beta$ sao cho MSE bé nhất bằng cách làm cho đạo hàm của MSE theo từng biến tiến gần đến 0 nhất có thể.
    \begin{center}
    $\frac{\partial L}{\partial X_i}=-\frac{2}{n}\sum^n_{i=1}(\hat{y}_i-y_i)\times X_i$
    \end{center}
    
    Tại mỗi thời điểm, $\frac{\partial L}{\partial X_i}$ cho biết hướng của MSE theo biến $X_i$, do đó ta cần hạn chế sự tăng theo hướng đó bằng cách cập nhật:
    \begin{center}
    $\omega_i = \omega_i - lr \times\frac{\partial L}{\partial X_i}$
    \end{center}
    Với lr là learning rate, thể hiện tỉ lệ hạn chế.
\end{itemize}

\subsubsection{Quá trình huấn luyện}

Quá trình huấn luyện và kiểm định được thực hiện với epoch = 1000, learning rate = 0.01, random seed = 42, test size = 0.3.

\textbf{Số liệu thực tế:}

Trên tập train: Mô hình dự đoán khá chính xác (với Index 666, chênh lệch 1857,05, khoảng 0,43\% so với giá trị thực. Với Index 256, chênh lệch 47020,97, khoảng 1,67\% so với giá trị thực).\\
Tuy nhiên với các mẫu có giá trị thực nhỏ (khoảng 200000), dự đoán có chênh lệch khá lớn (với Index 629, chênh lệch 148622,04, khoảng 70,77\% so với giá trị thực).

Trên tập test: Mô hình dự đoán không chính xác bằng, nhưng đa số chênh lệch đều khoảng dưới 30\% (Với Index 59, chênh lệch 90832,19, khoảng 20,18\% so với giá trị thực. Với Index 206, chênh lệch 1105514,28, khoảng 32.04\% so với giá trị thực).

Cần lưu ý, những giá trị cực đại thường sẽ bị dự đoán sai do trong quá trình học bị ảnh hưởng bởi việc xử lí outliner và giới hạn giá trị dự đoán bởi "clip".

\textbf{Đánh giá}:
\begin{itemize}
    \item Trên tập train, $R^2$ đạt 0.8968, MSE đạt 0.1035.
    \item Trên tập valid, $R^2$ đạt 0.8723, MSE đạt 0.1269.
    \item Trên Cross-Validation, $R^2$ đạt 0.8964, MSE đạt 0.0996.
\end{itemize}

Nhìn chung, khả năng dự đoán của mô hình trên thông số tương đối tốt, có thể thực hiện bài toán dự đoán giá xe. Tuy nhiên vẫn có thể cải thiện thêm hiệu suất huấn luyện bằng cách điền những giá trị còn thiếu ở những mẫu xe cùng dòng (Ví dụ với dòng A4 2.0 TDI (143 bhp) của hãng Audi, có thể điền những giá trị trống của Index 218 bằng giá trị của Index 220) trước khi điền bằng trung vị.
Một lưu ý cần quan tâm đến là mô hình sử dụng đặc trưng Make để huấn luyện, thuộc 1 trong 31 hãng xe trong dữ liệu. Trong trường hợp dữ đoán giá của những hãng xe khác, sẽ dẫn đến sai sót đáng kể, do đó cần xem xét việc xây dựng một mô hình tập trung vào các đặc trưng số thay vì các đặc trưng phân loại.